% --- Hauptebene: Kapitel ---
\chapter{Systemarchitektur und Dynamik}
\label{ch:architektur}

Dieses Kapitel beschreibt den grundlegenden Aufbau des Modells und stützt sich auf die theoretischen Konzepte der nichtlinearen Dynamik \cite{strogatz2018}.

% --- Zweite Ebene: Sektion ---
\section{Definition der Subsysteme}
\label{sec:subsysteme}

Das Gesamtsystem wird in drei gekoppelte Subsysteme unterteilt. Wie in Abbildung \ref{fig:regelkreis} zu sehen ist, spielt die Rückkopplung eine zentrale Rolle für die Stabilität.

% --- Einbindung eines Bildes mit Beschreibung ---
\begin{figure}[htbp]
    \centering
    % Ersetze 'regelkreis.png' durch deinen tatsächlichen Dateinamen
    % \includegraphics[width=0.8\textwidth]{images/regelkreis.png}
    
    % \caption[Kurztext fürs Verzeichnis]{Langtext für unter das Bild}
    \caption[Blockschaltbild des gekoppelten Regelkreises]{Detaillierte Darstellung der Rückkopplungsschleifen innerhalb des Hauptsystems. Beachten Sie die Verzögerungsglieder in der Feedback-Schleife, welche die Oszillationen im Systemzustand verursachen.}
    \label{fig:regelkreis}
\end{figure}

% --- Dritte Ebene: Untersektion ---
\subsection{Mathematische Formulierung der Kopplung}

Die Interaktion zwischen den Modulen lässt sich durch folgende Matrixgleichung beschreiben, wobei $K$ die Kopplungsmatrix darstellt:
$$ \dot{X} = F(X) + K \cdot G(X) $$
